% !TeX encoding = UTF-8
% !TeX spellcheck = pl_PL

\def\filename{rozdzial3}
\chapter{Propozycja systemu do diagnostyki zmian skórnych}

\section{Dostępne zbiory danych zmian skórnych}

Do rozwoju metod sztucznej inteligencji w dermatologii w istotny sposób przyczyniła się organizacja International Skin Imaging Collaboration (ISIC), której celem jest zwiększenie skuteczności wczesnego wykrywania czerniaka i tym samym zmniejszenie śmiertelności związanej z nowotworami skóry.
Od 2016 roku ISIC udostępnia obszerne zbiory obrazów z badań dermoskopowych wraz z danymi klinicznymi i wynikami diagnoz, organizując coroczne konkursy mające na celu doskonalenie algorytmów wizji komputerowej.
Archiwa ISIC odegrały kluczową rolę w rozwoju badań nad wykorzystaniem sztucznej inteligencji w diagnostyce nowotworów skóry, a opracowane przez organizację standardy są powszechnie akceptowane w dziedzinie CADx \cite{Li_Koban_Schenck_Giunta_Li_Sun_2022}


\subsection{Dobór zbiorów i opis zmian skórnych}

\section{Infrastruktura}

\section{Zastosowanie}

\section{Podsumowanie}